% ===== FILE 7/9: Eta (Nyquist), Ngưỡng & đa view, Split dị hướng (lưới), Split dị hướng (bậc thang k), Prune đa view =====

% --- Slide: Tỉ số eta — trực giác Nyquist ---
\begin{frame}{Bước 5: $\eta=\ell/w_{\min}$ --- trực giác Nyquist}
\scriptsize
Giờ đã có cả tử số ($\ell$, kích thước Gaussian đã chiếu) lẫn mẫu số ($w_{\min}$, bước sóng chi tiết cục bộ), định nghĩa
tỉ số cấu trúc:
\[
\boxed{\;\eta_j=\dfrac{\ell_j}{w_{\min}}\;}\qquad(\texttt{eta\_compute\_mode="wavelength"}\text{, mặc định}, \texttt{freq\_utils.py:348})
\]
Đây là một ngưỡng \textbf{hình học thuần tuý} (so sánh hai độ dài), \textbf{không phải} một mô hình vật lý về suy hao
biên độ tín hiệu --- tên gọi "trực giác Nyquist" chỉ mượn ý tưởng chung của lý thuyết lấy mẫu: muốn biểu diễn đúng một
dao động, "đơn vị lấy mẫu" (ở đây là Gaussian) phải đủ nhỏ so với chu kỳ của dao động đó.
\begin{itemize}\itemsep1pt
    \item $\eta>1$: $\ell$ \textbf{rộng hơn cả một chu kỳ} chi tiết $\Rightarrow$ một Gaussian không đủ khả năng bám
    theo dao động đó $\Rightarrow$ ứng viên \textbf{SPLIT}.
    \item $0.1<\eta\le1$: kích thước tương xứng với chi tiết cần biểu diễn $\Rightarrow$ \textbf{GIỮ NGUYÊN}.
    \item $\eta\le0.1$: $\ell$ nhỏ hơn $w_{\min}$ tới hơn $10$ lần --- Gaussian đang "phủ thừa" một vùng gần như không
    đổi $\Rightarrow$ ứng viên \textbf{PRUNE}.
    \item \textbf{Lưu ý về tài liệu nội bộ code:} một số dòng docstring (\texttt{gaussian\_model.py:837--840}) mô tả
    $\eta\sim(\sigma\omega)^2$ (quan hệ \emph{bậc hai}) --- nhưng công thức \textbf{thực sự được tính và dùng} ở trên là
    quan hệ \textbf{tuyến tính} theo kích thước. Toàn bộ bài giảng này bám theo công thức \emph{thực chạy}, không theo
    docstring (vốn có thể là tài liệu chưa cập nhật theo code).
\end{itemize}
\begin{center}
\includegraphics[width=0.62\linewidth]{02_nyquist.png}
\end{center}
\end{frame}

% --- Slide: Ngưỡng phân loại & tích luỹ đa view ---
\begin{frame}{Bước 6: ngưỡng phân loại và tích luỹ NHẤT QUÁN ĐA VIEW}
\scriptsize
Ngưỡng cố định trong code: \texttt{TAU\_HIGH=1.0}, \texttt{TAU\_LOW=0.1} (\texttt{freq\_utils.py:395--396}). Mỗi lần
render một view, lấy $\eta_{\max}=\max_j\eta_j$ (giá trị lớn nhất trong 3 trục) và cộng $1$ vào đúng một trong ba bộ đếm:
\[
\texttt{eta\_high\_count} \;(\eta_{\max}>1)\qquad
\texttt{eta\_mid\_count}\qquad
\texttt{eta\_low\_count}\;(\eta_{\max}\le0.1)
\]
Tới kỳ densify (mỗi $100$ iteration), các bộ đếm này được quy ra \textbf{tỉ lệ nhất quán qua nhiều view}
(\texttt{train.py:341--353}) --- đây chính là điểm khác biệt lớn nhất về \emph{triết lý} so với 3DGS gốc, vốn chỉ nhìn
vào gradient tích luỹ tại \textbf{một} thời điểm cuối cùng để quyết định, không phân biệt rõ ràng "có bao nhiêu góc
nhìn đồng ý với nhau":
\[
\text{high\_ratio}=\frac{\texttt{eta\_high\_count}}{\texttt{accum\_view\_count}}>0.8\qquad
\text{low\_ratio}=\frac{\texttt{eta\_low\_count}}{\texttt{accum\_view\_count}}>0.8
\]
\begin{itemize}\itemsep1pt
    \item \textbf{Trực giác:} chỉ khi \textbf{đa số áp đảo} ($>80\%$) các góc nhìn quan sát được đều đồng thuận rằng
    Gaussian "quá to" (hoặc "quá nhỏ"), quyết định split/prune mới được kích hoạt --- giúp tránh phản ứng thái quá theo
    nhiễu của một góc chụp bất thường (ví dụ do occlusion, ánh sáng loá, hay góc nhìn gần như cạnh của ellipsoid).
\end{itemize}
\begin{center}
\includegraphics[width=0.4\linewidth]{03_timeline_reset.png}
\end{center}
\end{frame}

% --- Slide: Split dị hướng — lưới chia tất định ---
\begin{frame}{Split dị hướng (1/2): chia đúng trục, đặt con trên lưới tất định}
\scriptsize
\texttt{densify\_and\_split\_structgs} (\texttt{gaussian\_model.py:638--}) --- phiên bản split \textbf{mới} của SADGS,
khi có sẵn \texttt{max\_eta\_3ch}, tính riêng cho từng trục:
\[
\boxed{\;k_{\text{axis}}=\big\lceil\sqrt{\max(\eta_{\text{axis}},1)}\big\rceil\;},\qquad
N_{\text{con}}=k_x\,k_y\,k_z,\qquad s_{\text{con}}=\frac{s_{\text{cha}}}{k_{\text{axis}}^{\,p}}\ (p=\texttt{ks\_scale\_power}=1.0)
\]
\begin{itemize}\itemsep1pt
    \item Điểm khác biệt lớn với split nguyên bản (slide trước, phần 1): $k$ được tính \textbf{riêng cho từng trục}
    --- trục nào thực sự vi phạm ($\eta_{\text{axis}}>1$) mới bị chia; trục còn "khoẻ mạnh" giữ nguyên $k=1$.
    \item Vị trí các con được đặt trên một \textbf{lưới tất định} (deterministic grid) --- \textbf{không} lấy mẫu ngẫu
    nhiên như cơ chế gốc: khoảng cách tâm-tâm giữa các con liền kề $=\sigma_{\text{con}}\sqrt{12}$, toàn bộ lưới được
    xoay theo đúng hướng $R$ của Gaussian cha (\texttt{gaussian\_model.py:766--780}).
    \item Có thể kiểm chứng lưới này \textbf{bảo toàn phương sai}: tổng phương sai của các con đặt cách đều theo công
    thức trên đúng bằng phương sai của cha ban đầu --- tức đây không phải một lựa chọn tuỳ ý mà có cơ sở thống kê.
\end{itemize}
\begin{center}
\includegraphics[width=0.4\linewidth]{04_split_grid_2d.png}
\end{center}
\end{frame}

% --- Slide: Split dị hướng — bậc thang k theo eta ---
\begin{frame}{Split dị hướng (2/2): $k$ tăng theo bậc thang khi $\eta$ tăng}
\scriptsize
Vì $k_{\text{axis}}=\lceil\sqrt{\max(\eta_{\text{axis}},1)}\rceil$ là một hàm \textbf{làm tròn lên}, quan hệ giữa $k$ và
$\eta$ không mượt mà là một \textbf{dạng bậc thang}: $k$ chỉ nhảy lên một mức mới khi $\eta$ vượt qua các ngưỡng $1,4,9,16,\dots$
($k^2$). Điều này có nghĩa:
\begin{itemize}\itemsep1pt
    \item Với $\eta$ vừa mới vượt $1$ (ví dụ $\eta=1.2$): $k=2$, sinh $N_{\text{con}}=2$ trên trục đó --- một mức tăng
    "thận trọng", tương xứng với mức độ vi phạm còn nhẹ.
    \item Với $\eta$ lớn hơn nhiều (ví dụ $\eta=20$): $k=\lceil\sqrt{20}\rceil=5$, sinh hẳn $5$ con trên trục đó --- phản
    ứng \textbf{mạnh tay hơn} tương xứng với mức độ vi phạm nghiêm trọng hơn.
    \item So với split nguyên bản (luôn cố định $N=2$ bất kể mức độ "quá to" ra sao), cách làm bậc thang này cho phép hệ
    thống \textbf{tự điều chỉnh cường độ phản ứng} theo đúng quy mô của vấn đề --- tránh vừa chia không đủ (còn dư
    $\eta>1$ sau một lần split) vừa tránh chia quá tay gây bùng nổ số lượng Gaussian không cần thiết.
\end{itemize}
\begin{center}
\includegraphics[width=0.4\linewidth]{04_k_eta_stair.png}
\end{center}
\end{frame}

% --- Slide: Prune theo nhất quán đa view ---
\begin{frame}{Prune theo cấu trúc: OR với tiêu chí gốc, không thay thế}
\scriptsize
Trong \texttt{densify\_and\_prune\_structgs} (\texttt{gaussian\_model.py:1009--1019}), mặt nạ prune \textbf{cấu trúc}
(low\_ratio$>0.8$, slide trước) được kết hợp bằng phép \textbf{OR} với các tiêu chí nguyên bản đã học ở phần ADC gốc
--- một Gaussian bị xoá nếu \textbf{bất kỳ} điều kiện nào sau đúng:
\[
\texttt{prune\_mask} = \underbrace{(\alpha<0.1)}_{\text{opacity thấp}} \;\lor\; \underbrace{(r_{2D}>\texttt{max\_screen\_size})}_{\text{quá to trên màn hình}} \;\lor\; \underbrace{(\max(s)>0.1\cdot\texttt{extent})}_{\text{quá to trong không gian}} \;\lor\; \underbrace{(\text{low\_ratio}>0.8)}_{\text{cấu trúc, đa view}}
\]
\begin{itemize}\itemsep1pt
    \item \texttt{min\_opacity}$=0.1$ trong nhánh này --- \textbf{cao hơn nhiều} so với $0.005$ ở nhánh 3DGS gốc, tức
    prune mạnh tay hơn theo tiêu chí opacity đơn thuần, không chỉ nhờ thêm điều kiện cấu trúc.
    \item Cả clone lẫn split (2 slide trước) còn bị chặn thêm bởi \texttt{metric\_mask}$=(\texttt{accum\_view\_count}>0.5)$
    --- một Gaussian vừa mới sinh, chưa từng được quan sát trong đợt tích luỹ hiện tại, sẽ \textbf{không} bị động tới ở
    đợt densify đó, tránh phản ứng vội vàng dựa trên dữ liệu quá ít ỏi.
\end{itemize}
\begin{center}
\includegraphics[width=0.4\linewidth]{07_mat_na_prune_or.png}
\end{center}
\end{frame}
